% =====================================================================
\chapter{Background and Related Work}\label{ch:background}
% =====================================================================

% \begin{plan}[Chapter plan --- Background and related work]
% \pages{12 pages} (the largest chapter, and where most of the bibliography is
% used). \textbf{Job:} give the reader exactly the theory needed for
% Chapters 3--6 and end with a gap that the thesis fills. Nine sections, each
% ending with one sentence on ``what this means for the audit''. Three
% equations (attention, cross-attention, next-token softmax), one architecture
% figure, two tables (SIGIR queries, gap comparison).
% \end{plan}

% ---------------------------------------------------------------------
\section{Vision language models}\label{sec:bg-vlm}

The five models evaluated here belong to two families of decoder-based
vision language models. Both combine an image encoder with a language model,
but differ in how the visual features enter the decoder. Attention provides
the underlying operation, weighting values $V$ according to the similarity
between queries $Q$ and keys $K$ \cite{vaswani2017attention}:

\begin{equation}
  \operatorname{Attn}(Q,K,V) = \softmax\!\left(\frac{QK^{\top}}{\sqrt{d_k}}\right)V
  \label{eq:attention}
\end{equation}

A Vision Transformer represents an image through patch embeddings processed
by transformer layers \cite{dosovitskiy2021vit}. Contrastive training on
paired images and text can align their representations, as demonstrated
by CLIP \cite{radford2021clip}. These components provide a basis for
multimodal architectures, although their integration and training differ.

Llama-3.2-Vision connects image features to the language decoder through
gated cross-attention blocks \cite{grattafiori2024llama3}, following the
integration approach introduced by Flamingo \cite{alayrac2022flamingo}:

\begin{equation}
  \bm h_\ell' = \bm h_\ell + \tanh(\alpha_\ell)\,
  \operatorname{Attn}\!\big(\bm h_\ell W_Q,\; \bm e(v) W_K,\; \bm e(v) W_V\big)
  \label{eq:crossattn}
\end{equation}

Here $\bm e(v)$ denotes image features and $\alpha_\ell$ is a learned gate
controlling the cross-attention contribution. The reference implementation
includes gates on the attention and feedforward sublayers
\cite{wolf2020transformers}. This permits visual features to influence
decoder states without inserting them as ordinary text tokens.

Qwen2-VL instead merges image patch representations into visual tokens
that enter the decoder sequence alongside text \cite{wang2024qwen2vl}.
Their number varies with image resolution, so visual input affects sequence
length as well as content. The experimental settings bound this token count
as described in \cref{sec:design-impl}; they do not make every input
sequence the same length. \Cref{fig:fusion} illustrates the two routes.

For the evaluated architectures, the final hidden state produces next-token
logits and their normalised probability distribution:

\begin{equation}
  p(t \mid x, v) = \softmax(\bm z)_t = \frac{\exp(z_t)}{\sum_{u\in\vocab}\exp(z_u)},
  \qquad \bm z = W_{\text{out}}\,\RMSNorm(\bm h_L)
  \label{eq:nexttoken}
\end{equation}

The audit reads this distribution before sampling, preserving probabilities
that would be reduced to a selected token during generation.
\Cref{ch:framework} defines the measurements obtained from that distribution.

\begin{figure}[htbp]\centering
\resizebox{\linewidth}{!}{%
\begin{tikzpicture}[node distance=5mm and 7mm,
  box/.style={draw,rounded corners,align=center,font=\footnotesize,minimum
height=7mm,inner sep=3pt},
  tok/.style={draw,align=center,font=\scriptsize,minimum height=5.5mm,inner sep=2.5pt},
  arr/.style={-{Stealth[length=2mm]},semithick},
  hdr/.style={font=\footnotesize\itshape}]

% ---- left panel: cross-attention adapters -------------------------------
\node[hdr] (lh) {Llama-3.2-Vision: gated cross attention};
\node[box,below=4mm of lh.south west,anchor=north west] (lvit) {ViT-H/14\\\scriptsize 32 layers, width 1280};
\node[box,below=of lvit] (lca) {cross-attention blocks\\\scriptsize layers 3, 8, 13, \ldots, 38};
\node[box,below=of lca] (lsa) {self-attention layers\\\scriptsize unchanged};
\node[box,below=of lsa,fill=blue!8] (lhead) {LM head};
\draw[arr] (lvit) -- node[right,font=\scriptsize] {$\bm e(v)$} (lca);
\draw[arr] (lca) -- (lsa);
\draw[arr] (lsa) -- (lhead);
\node[left=5mm of lsa,font=\scriptsize,align=right] (ltext) {text\\tokens};
\draw[arr] (ltext) -- (lsa);
\node[left=5mm of lvit,font=\scriptsize] (lim) {image $v$};
\draw[arr] (lim) -- (lvit);

% ---- right panel: visual tokens in the sequence --------------------------
\node[hdr,right=34mm of lh.east,anchor=west] (rh) {Qwen2-VL: visual tokens in the sequence};
\node[box,below=4mm of rh.south west,anchor=north west] (rvit) {ViT at native resolution};
\node[box,below=of rvit] (rmrg) {patch merger};
\node[tok,below=of rmrg] (rseq) {$[\,$visual tokens $\;\Vert\;$ text tokens$\,]$};
\node[box,below=of rseq] (rdec) {decoder layers\\\scriptsize unchanged};
\node[box,below=of rdec,fill=blue!8] (rhead) {LM head};
\draw[arr] (rvit) -- (rmrg);
\draw[arr] (rmrg) -- (rseq);
\draw[arr] (rseq) -- (rdec);
\draw[arr] (rdec) -- (rhead);
\node[left=5mm of rvit,font=\scriptsize] (rim) {image $v$};
\draw[arr] (rim) -- (rvit);

% ---- shared output -------------------------------------------------------
\node[box,below=9mm of $(lhead.south)!0.5!(rhead.south)$,fill=blue!8]
  (dist) {next-token distribution $p(\cdot\mid x,v)$};
\draw[arr] (lhead) |- ([yshift=4mm]dist.north) ;
\draw[arr] (rhead) |- ([yshift=4mm]dist.north);
\draw[arr] ([yshift=4mm]dist.north) -- (dist);
\node[below=2.5mm of dist,font=\footnotesize,text=blue!60!black]
  {\textbf{measured here} (\cref{eq:nexttoken})};
\end{tikzpicture}}
\caption{Two ways of injecting an image into a language model.
Llama-3.2-Vision uses gated cross-attention blocks to connect image features
to the decoder; the illustrated Llama-11B configuration places eight such
blocks among 40 decoder layers. Qwen2-VL emits visual tokens that join the
text tokens in the input sequence, so the sequence length grows with the
image. Both routes end at the same distribution, which is where this thesis
takes its reading.}
\label{fig:fusion}
\end{figure}

% \begin{plan}
% \begin{itemize}
%   \item Transformer attention \cite{vaswani2017attention} --- give
%         \cref{eq:attention}; vision transformers \cite{dosovitskiy2021vit};
%         contrastive image--text alignment \cite{radford2021clip}.
%   \item Two ways of fusing vision into a language model, both audited here:
%         (a) \emph{cross-attention adapters} (Flamingo-style
%         \cite{alayrac2022flamingo}; Llama~3.2-Vision inserts cross-attention
%         layers after every fourth self-attention layer of the language model,
%         fed by a ViT-H/14 encoder \cite{grattafiori2024llama3}) --- give
%         \cref{eq:crossattn}; (b) \emph{visual tokens in the input sequence}
%         with dynamic resolution (Qwen2-VL \cite{wang2024qwen2vl}).
%   \item Why it matters for the audit: in (a) the image enters through
%         dedicated layers, so a grey patch can change behaviour without
%         changing the prompt; in (b) image size changes the sequence length,
%         which is why the visual token budget is pinned (\cref{sec:design-impl}).
%   \item End with the output layer: \cref{eq:nexttoken}, the quantity this
%         thesis reads.
% \end{itemize}
% \end{plan}
% \begin{plan}[Notes on the equations above]
% \cref{eq:crossattn} is the gated cross-attention form of
% \cite{alayrac2022flamingo}; $\bm e(v)$ are image-encoder features and
% $\alpha_\ell$ a learned gate. \textbf{Checked (Nov 2026):} the Llama-3.2
% adapter \emph{is} tanh-gated. \texttt{MllamaCrossAttentionDecoderLayer} in the
% reference implementation holds \texttt{cross\_attn\_attn\_gate} and
% \texttt{cross\_attn\_mlp\_gate}, both \texttt{nn.Parameter(torch.zeros(1))},
% applied as \texttt{residual + gate.tanh() * hidden\_states}; default
% \texttt{cross\_attention\_layers} are $[3,8,13,18,23,28,33,38]$ and the vision
% config is patch 14 / width 1280 / 32 layers / image 448. So \cref{eq:crossattn}
% may be presented as the actual form, not only ``the general form''.
% \cref{eq:nexttoken} reappears in \cref{prop:uniform} (why a zero hidden state
% gives a uniform distribution).
% \end{plan}
% ---------------------------------------------------------------------

\section{Social bias in learned representations}\label{sec:bg-bias}

Learned representations can reproduce social associations present in their
training material. Caliskan et al.\ \cite{caliskan2017semantics} adapted
the Implicit Association Test \cite{greenwald1998iat} to word embeddings,
recovering associations involving race, gender and occupations from a web
corpus. Such results demonstrate measurable associations in representation
geometry, without requiring a model to state an explicit belief.

Visual systems also show differences across demographic groups.
Buolamwini and Gebru \cite{buolamwini2018gender} found that commercial
gender-classification error rates varied across skin tone and gender,
with the largest errors for darker-skinned women. In multimodal evaluation,
VisoGender examines gendered reference resolution in occupational images
\cite{hall2023visogender}, while ModSCAN measures stereotyped responses
across occupations, descriptors and persona traits
\cite{jiang2024modscan}. These studies examine different tasks and
measures, making their scope important when comparing findings.

Component-level analysis provides evidence about where particular effects
arise. Weng et al.\ \cite{weng2024images} used causal mediation in an
object detection model and found image features contributed more to their
bias measures than text features or the fusion encoder on the two evaluated
datasets. This motivates examining visual input in an audit, but does not
identify the responsible components in the generative models studied here.

Associations measured within representations can also relate to downstream
retrieval. Ghate et al.\ \cite{ghate2025biases} reported a mean correlation
of $\rho=0.83\pm0.10$ between intrinsic and extrinsic bias measures across
114 analyses, six social groups and three encoder-based models.
Their findings motivate testing downstream behaviour alongside intrinsic
measurements, rather than assuming that the relationship transfers to a
different model or task.

Interpretation also requires specifying what a bias measure represents.
Blodgett et al.\ \cite{blodgett2020language} emphasise the need to identify
the relevant construct and potential harm, while Gallegos et al.\
\cite{gallegos2024bias} distinguish measurements at the embedding,
probability and generated-text levels. The present study examines
probabilities and gender-related trait attribution, with the limits of the
adapted construct stated separately from possible consequences in deployment.

% \begin{plan}
% \begin{itemize}
%   \item Human-like associations in embeddings, measured with an IAT
%         analogue \cite{caliskan2017semantics,greenwald1998iat}.
%   \item Vision: intersectional error disparities \cite{buolamwini2018gender};
%         VLM benchmarks for gender in image--text tasks
%         \cite{hall2023visogender}; stereotype bias from both modalities
%         \cite{jiang2024modscan}; image features as the main mediator
%         \cite{weng2024images}; propagation to retrieval \cite{ghate2025biases}.
%   \item What ``bias'' means and how it is measured: \cite{blodgett2020language}
%         (be explicit about the harm and the construct);
%         \cite{gallegos2024bias} classifies metrics as embedding-,
%         probability- or generation-based --- place this thesis in the
%         \emph{probability} level.
% \end{itemize}
% \end{plan}
% ---------------------------------------------------------------------

\section{Gender bias in image search and its human perception}\label{sec:bg-search}

Image search can both reflect social stereotypes and influence users'
perceptions. Kay et al.\ \cite{kay2015unequal} found occupational gender
imbalances relative to labour statistics and showed that exposure to results
shifted participants' estimates of the occupational distribution.
Metaxa et al.\ \cite{metaxa2021image} subsequently examined gender and
race representation, identifying continuing distortions in search results.

The pattern extends beyond occupations to descriptions of personal traits.
Otterbacher et al.\ \cite{otterbacher2017competent} found women associated
with warm traits and men with agentic traits in image search.
This corresponds with the warmth and competence dimensions of the Stereotype
Content Model \cite{fiske2002scm}, which distinguishes the content of
stereotypes from representation counts alone.

The 2018 study provides the human comparison used in this thesis
\cite{otterbacher2018sigir}. Otterbacher, Checco, Demartini and Clough
recruited 281 crowd workers across the United Kingdom, United States and
India, using nine-image grids for ten queries. Their reverse image search
procedure first asked participants to infer the query from a grid, before
revealing the actual query and requesting an objectivity judgement.
Ratings ranged from 1, extremely subjective, to 7, extremely objective;
the ASI was administered afterwards to reduce priming from its questions.
Participants therefore described the visual content before being asked to
judge the search engine's objectivity, rather than beginning with an explicit
gender-bias question. Although the released dataset contains ten queries,
the published initial analysis focused on four gender-skewed trait queries
and the balloon control. One participant with undisclosed gender was excluded
from that analysis, accounting for the reported sample of 280.

Controlling for country, age, gender and self-reported search experience,
higher benevolent sexism was associated with a lower likelihood of judging
the ``smart person'' and ``warm person'' grids as non-objective.
Those grids were predominantly male and female, respectively.
The association did not appear for ``aggressive person'' or ``anxious
person''. For the balloon control, 28 of 280 participants judged the
results non-objective, with no measured participant characteristic predicting
that judgement.

The association was specific to benevolent sexism and these queries.
It motivates examining whether gender-related inventory responses in models
correspond with their judgements of image search results, with the original
human findings providing the starting point for that comparison.

\begin{table}[htbp]\centering\small
\begin{tabular}{@{}lccc@{}}
\toprule
Query (as in the dataset) & Trait valence & Grid skew & Low-ASI mean rating \\
\midrule
smart person       & $+$ & M  & 5.257 \\
aggressive person  & $-$ & M  & 5.486 \\
warm person        & $+$ & F  & 4.899 \\
anxious person     & $-$ & F  & 5.453 \\
hot air baloon     & $=$ & n/a & 6.365 \\
bossy person       & \multicolumn{2}{c}{\datagap{not reported}} & 5.020 \\
calm person        & \multicolumn{2}{c}{\datagap{not reported}} & 5.473 \\
casual person      & \multicolumn{2}{c}{\datagap{not reported}} & 4.703 \\
interested person  & \multicolumn{2}{c}{\datagap{not reported}} & 4.804 \\
working person     & \multicolumn{2}{c}{\datagap{not reported}} & 5.338 \\
\bottomrule
\end{tabular}
\caption{The ten SIGIR 2018 queries. Valence and skew for the first five are
from Table~1 of \cite{otterbacher2018sigir}, which lists those five only;
the valence and skew of the remaining five are not reported there and cannot
be recovered from the released responses, which carry ratings rather than
image annotations. The low-ASI mean is the per-query baseline $\rlow$ used in
\cref{ch:study2} (median split on the participants' total ASI score). ``baloon'' is the dataset's spelling and was used
verbatim in the prompts.}
\label{tab:sigir-queries}
\end{table}

\Cref{tab:sigir-queries} lists the queries and baseline ratings used here;
\cref{sec:design-phase2} explains how the human groups and image task are
adapted for the model comparison.

% \begin{plan}
% \begin{itemize}
%   \item Under-representation and stereotype exaggeration for occupations,
%         and its effect on users' beliefs \cite{kay2015unequal}; little change
%         five years later \cite{metaxa2021image}.
%   \item Character traits rather than occupations: women shown for warm
%         traits, men for agentic ones \cite{otterbacher2017competent}, grounded
%         in the warmth/competence model \cite{fiske2002scm}.
%   \item \textbf{The human baseline used in this thesis}
%         \cite{otterbacher2018sigir}: crowdsourced, three countries (UK, USA,
%         India), 281 participants, 10 queries, 3$\times$3 Google image grids,
%         objectivity rated 1 (extremely subjective) to 7 (extremely objective),
%         ASI administered at the end. Key result: higher benevolent sexism
%         predicts \emph{not} rating the ``smart person'' and ``warm person''
%         grids as non-objective --- sexist participants are less likely to see
%         bias. This sign matters for interpreting $\dm$
%         (\cref{sec:framework-extrinsic}).
%   \item \cref{tab:sigir-queries}: the 10 queries; fill trait valence and
%         bias direction from Table 1 of \cite{otterbacher2018sigir} for the
%         five queries it lists, and mark the other five as reported in the
%         released dataset.
% \end{itemize}
% \end{plan}
% ---------------------------------------------------------------------

\section{Ambivalent sexism and its measurement}\label{sec:bg-asi}

The Ambivalent Sexism Inventory (ASI) was developed to measure two related
components of attitudes towards women \cite{glick1996asi}.
Hostile sexism concerns antagonistic beliefs, including claims that women
exaggerate discrimination or seek control over men. Benevolent sexism uses
apparently favourable descriptions while assigning women restrictive roles,
such as portraying them as pure, fragile or dependent on male protection.
The inventory contains 22 agreement items, with eleven in each subscale.
Cross-national findings show that the subscales are positively related and
associated with measures of gender inequality \cite{glick2000beyond}.
The distinction matters because an approving tone can still assign different
expectations to women and men; benevolence in this framework is not simply
courtesy or a favourable evaluation.

The adapted inventory in \cref{tab:items} illustrates both components:
``exaggerating problems they have at work'' belongs to the hostile subscale,
while ``someone who should be cherished and protected'' belongs to the
benevolent subscale. These are the phrases administered to the models,
rather than verbatim statements from the original questionnaire.

Six items are reverse-keyed, meaning that agreement indicates lower
endorsement of the corresponding sexist belief; the remaining sixteen are
positively keyed. In the administered ordering in \cref{tab:items}, the
reverse-keyed items are 6, 7, 9, 15, 20 and 22.
The original human questionnaire and the signed probability index used here
are different scoring procedures, so the effect of keying must be examined
for the particular model measure. \Cref{prop:shift} establishes that effect
for the adapted index.

Response style provides a further reason to examine agreement separately
from item content. Acquiescence and preferences for particular scale positions
can distort questionnaire comparisons, with their effects depending on
instrument design \cite{baumgartner2001response}. Reverse-keyed items can
also introduce error when respondents process their wording differently
\cite{weijters2013reversed}. These concerns apply to human measurement
and motivate testing the assumptions again for model responses.

Validity concerns the inferences supported by scores rather than an inherent
property of the questionnaire \cite{messick1995validity}.
The ASI's human validation therefore does not automatically validate an
adaptation that changes the referent, response format and respondent.
\Cref{sec:design-phase1} describes those changes, while
\cref{sec:disc-threats} discusses their consequences for interpretation.

% \begin{plan}
% \begin{itemize}
%   \item Hostile vs benevolent sexism, the 22-item ASI, 11 items each
%         \cite{glick1996asi}; cross-cultural validity and the HS--BS
%         correlation \cite{glick2000beyond}. (Reserve:
%         \cite{glick2001alliance} on BS as a justification of inequality;
%         \cite{dardenne2007insidious} on BS harming performance --- use if you
%         want to argue BS is not ``harmless''.)
%   \item Keying: 16 items are scored so that agreement means more sexism and
%         6 are reverse-keyed (items 6, 7, 9, 15, 20, 22). Introduce this here;
%         \cref{prop:shift} shows why it matters for a model.
%   \item Response styles: acquiescence and its effect on scale scores depend on
%         the share of reversed items \cite{baumgartner2001response}; reversed
%         items carry their own bias \cite{weijters2013reversed}.
%   \item Validity as an argument, not a property \cite{messick1995validity}
%         --- sets up the threats-to-validity structure of \cref{sec:disc-threats}.
% \end{itemize}
% \end{plan}
% ---------------------------------------------------------------------

\section{Administering questionnaires to language models}\label{sec:bg-surveys}

Questionnaires offer a structured way to compare model responses, but their
use does not establish that a model behaves as a human respondent.
Dominguez-Olmedo et al.\ \cite{dominguez2024questioning} administered
American Community Survey questions to 43 models and identified strong
sensitivity to option position and labels, including preference for ``A''.
After randomising presentation, responses tended towards uniformity rather
than revealing a stable demographic profile.

Tjuatja et al.\ \cite{tjuatja2024response} examined whether nine models
reproduced response biases documented in human surveys. The models generally
differed from human respondents, with reinforcement-learning-from-human-feedback
models showing the largest divergence. They also changed responses under
perturbations that did not similarly affect human answers.
These results caution against interpreting equivalent questionnaire formats
as evidence of equivalent response processes.

Probability-based surveys raise related questions about whose opinions are
represented. Santurkar et al.\ \cite{santurkar2023opinions} found that
recovered opinion distributions did not closely match any evaluated American
demographic group. Sharma et al.\ \cite{sharma2024sycophancy} examined
agreement with user views and linked sycophantic responses to human preference
signals in training. Agreement can therefore reflect influences other than
the content of the proposition being evaluated.

For the present audit, these findings motivate checking prompt sensitivity
and separating general agreement from discrimination between item types.
The decomposition in \cref{sec:framework-responsebias} addresses the specific
sensitivity of the adapted inventory score, rather than assuming every
questionnaire or model is affected in the same way.

% \begin{plan}
% \begin{itemize}
%   \item Models show ordering/labelling biases on survey items, and once these
%         are controlled their answers tend toward uniform
%         \cite{dominguez2024questioning}.
%   \item RLHF-tuned models do not reproduce human response biases and react to
%         perturbations humans ignore \cite{tjuatja2024response}.
%   \item (Reserve) opinion alignment from answer-option probabilities
%         \cite{santurkar2023opinions}; agreement with the user
%         (sycophancy) as an RLHF effect \cite{sharma2024sycophancy} --- the
%         closest model analogue of acquiescence.
%   \item What this means here: a model's yes-rate is a free parameter, so an
%         inventory score must be read together with its response style.
% \end{itemize}
% \end{plan}
% ---------------------------------------------------------------------

\section{Alignment, refusal and covert bias}\label{sec:bg-alignment}

Instruction tuning and alignment modify the model that produces the response
distribution, rather than acting only after probabilities have been calculated.
Training on human preferences \cite{ouyang2022training} or written principles
\cite{bai2022constitutional} can consequently affect both explicit answers
and probability-based measurements. Differences between explicit responses
and indirectly measured associations require empirical investigation.

R\"ottger et al.\ \cite{rottger2024xstest} documented exaggerated safety,
where models refuse harmless prompts resembling unsafe requests.
Arditi et al.\ \cite{arditi2024refusal} identified a direction in the residual
stream whose manipulation could suppress or induce refusal in the evaluated
models. This provides evidence for an internal mechanism, but does not make
low probability on a selected response set a universal refusal detector.

Explicitly favourable responses can also coexist with negative indirect
associations. Hofmann et al.\ \cite{hofmann2024covert} compared overt racial
stereotypes with judgements associated with African American English,
including simulated employment and sentencing decisions. Human preference
training reduced overt expression without correspondingly removing the covert
associations. Bai et al.\ \cite{bai2025explicitly} similarly identified
stereotype associations in models that passed explicit bias tests.
These findings concern the evaluated measures and do not imply that every
apparent difference reveals a concealed attitude.

Measurement framing introduces another source of variation. Gao and Kreiss
\cite{gao2025measuring} found that gender-bias estimates changed with
evaluative framing and that discrete-choice measures amplified apparent bias
relative to probability measures. Persona framing also affects social
reasoning \cite{yang2026persona}. Such sensitivity motivates the hypothesis
of algorithmic defensiveness introduced in \cref{sec:intro-problem}, but
changes in response distributions alone do not demonstrate that a model
recognises an audit or deliberately conceals its associations.

% \begin{plan}
% \begin{itemize}
%   \item RLHF and constitutional training \cite{ouyang2022training,bai2022constitutional}
%         (reserve for multimodal RLHF: \cite{sun2024aligning}).
%   \item Over-refusal of safe prompts that resemble unsafe ones
%         \cite{rottger2024xstest} --- the closest text analogue of Llama
%         declining a yes/no question about ``this person'' in a grey square.
%   \item Refusal is represented internally (a single direction in activations)
%         \cite{arditi2024refusal} --- supports measuring it before decoding.
%   \item Overt vs covert bias \cite{hofmann2024covert,bai2025explicitly};
%         framing effects \cite{gao2025measuring,yang2026persona}. Note that
%         \cite{gao2025measuring} also report that token-probability measures
%         and discrete-choice measures give different bias estimates --- an
%         argument for reporting the measurement level explicitly.
% \end{itemize}
% \end{plan}
% ---------------------------------------------------------------------

\section{Evaluation from probabilities}\label{sec:bg-probability}

Probability measurements require a response set that represents the intended
answers. Holtzman et al.\ \cite{holtzman2021surface} showed that competing
surface forms can divide probability among alternative expressions of the
same answer. The pooled measure in \cref{eq:pooled} accounts for specified
case variants, while \cref{sec:design-faults} records the consequences of
using an incomplete response set.

Answer-token priors introduce a separate influence. Zhao et al.\
\cite{zhao2021calibrate} used content-free inputs to estimate preferences
that arise without task evidence and improved classification through
calibration. This motivates distinguishing answer preferences from responses
to item content; the visual controls in \cref{sec:design-stimuli} examine
responses when informative image content is absent, without implementing
that text-calibration procedure directly.

The suitability of probability measurements also depends on the question
being evaluated. Hu and Levy \cite{hu2023prompting} found that probability
measurements outperformed prompted metalinguistic judgements in their tasks.
Wang et al.\ \cite{wang2024myanswer}, however, found substantial mismatches
between first-token preferences and generated answers on multiple-choice
questions, with reported mismatch rates above 60\% in their evaluation.
Even high probability coverage at the first response position does not
therefore guarantee agreement with a subsequent generated answer. These studies examine different evaluation settings and show why
probabilities and generated text should not be treated as interchangeable
measurements. Captured mass reports response coverage but does not resolve
that distinction.

Preferences can attach to option identifiers as well as answer content.
Zheng et al.\ \cite{zheng2024robust} examined 20 models on three benchmarks
and found sensitivity to the position and identifiers of answer options.
The rating task here also uses discrete labels, motivating a reversal of
the scale endpoints to test whether normalised judgements remain consistent.

% \begin{plan}
% \begin{itemize}
%   \item Probability mass is split across surface forms of the same answer
%         \cite{holtzman2021surface} --- the theoretical reason for pooling
%         \{yes, Yes, YES\} (\cref{eq:pooled}).
%   \item Models carry answer priors that content-free inputs reveal
%         \cite{zhao2021calibrate} --- the noise and grey-patch stimuli are
%         content-free inputs in exactly this sense.
%   \item Direct probability measurements outperform prompting the model for a
%         judgement \cite{hu2023prompting}; but first-token probabilities and
%         text answers can disagree, especially for safety-tuned models
%         \cite{wang2024myanswer} --- state this as a limitation, not a
%         refutation.
%   \item Option-ID token priors cause selection bias \cite{zheng2024robust}
%         --- the digit prior seen in \cref{sec:s2-deviation}.
% \end{itemize}
% \end{plan}
% ---------------------------------------------------------------------

\section{LLM and MLLM judges, and positional perturbation}\label{sec:bg-judges}

Automated judges can reproduce systematic preferences unrelated to the
quality they are intended to evaluate. Zheng et al.\ \cite{zheng2023judging}
identified position bias, verbosity bias and preference for a judge's own
outputs. Wang et al.\ \cite{wang2024fair} showed that changing candidate
order could reverse rankings, demonstrating the need to control presentation
when interpreting model judgements.

Chen et al.\ \cite{chen2024mllmjudge} evaluated multimodal judges on scoring,
pairwise comparison and batch ranking. Agreement with human preferences
varied by task, with scoring and batch ranking less reliable than pairwise
comparison in their evaluation. Their analysis also identified position,
length and egocentric biases, making performance on one judgement format
insufficient evidence of reliability on another.

CALM evaluates twelve judgement biases through controlled input perturbations
that should preserve the verdict \cite{ye2025justice}.
This approach tests the stability of a judgement separately from its
agreement with a reference answer. For example, exchanging the order of two
candidates changes their presentation without changing their substantive
content, allowing the resulting preference to be checked for order sensitivity.
The present study adapts its robustness-rate concept to scale reversal,
checking whether the same normalised rating is retained after the endpoints
are exchanged. The expected raw digit changes under reversal, but its
interpretation should remain the same after normalisation; this distinguishes
scale adherence from a preference for particular numerical labels.
Related work examines the ordering of scoring rubrics and
identifiers \cite{li2026scoring} and compares model judgement biases with
human biases \cite{chen2024humans}. These findings motivate the robustness
check without establishing that instability in the present task has the
same cause.

% \begin{plan}
% \begin{itemize}
%   \item Position, verbosity and self-enhancement biases
%         \cite{zheng2023judging}; order of candidates changes rankings
%         \cite{wang2024fair}; multimodal judges diverge from humans on scoring
%         \cite{chen2024mllmjudge}.
%   \item CALM: a perturb-and-compare framework with robustness and consistency
%         rates \cite{ye2025justice} --- the source of $\RR$ (\cref{eq:rr}).
%   \item (Reserve) rubric-order and score-ID biases in scoring judges
%         \cite{li2026scoring}; judgement biases of humans vs LLMs
%         \cite{chen2024humans}.
% \end{itemize}
% \end{plan}
% ---------------------------------------------------------------------

\section{Research gap}\label{sec:bg-gap}

\begin{table}[htbp]\centering\small
\begin{tabularx}{\linewidth}{@{}lc>{\raggedright\arraybackslash}X>{\raggedright\arraybackslash}Xcc@{}}
\toprule
Work & Images & Level measured & Human anchor & \makecell{Order\\control} & \makecell{Coverage\\statistic} \\
\midrule
\cite{caliskan2017semantics} & -- & embeddings & IAT results & -- & -- \\
\cite{ghate2025biases}       & \checkmark & embeddings $\to$ retrieval & -- & -- & -- \\
\cite{jiang2024modscan}      & \checkmark & outputs & -- & -- & -- \\
\cite{weng2024images}        & \checkmark & internal (mediation) & -- & -- & -- \\
\cite{hofmann2024covert}     & -- & probabilities & human stereotype studies & -- & -- \\
\cite{bai2025explicitly}     & -- & outputs & IAT-style & -- & -- \\
\cite{gao2025measuring}      & -- & outputs + probabilities & -- & -- & -- \\
\cite{ye2025justice}         & -- & judge outputs & -- & \checkmark & -- \\
\cite{chen2024mllmjudge}     & \checkmark & judge outputs & human agreement & \checkmark & -- \\
\cite{otterbacher2018sigir}  & \checkmark & human ratings & ASI & -- & -- \\
\midrule
\textbf{This thesis}         & \checkmark & probabilities & ASI + SIGIR 2018 & \checkmark & $\capm$ \\
\bottomrule
\end{tabularx}
\caption{Positioning against prior work. Each row records what the paper
reports in its abstract and described method; a dash means the element is not
reported. Checked against each source in November 2026.}
\label{tab:gap}
\end{table}

The approaches compared in \cref{tab:gap} examine social associations,
response sensitivity, human perceptions and judgement robustness, but do not
provide an audit combining visual inputs, targeted next-token probability
measurements, explicit response-set coverage and the human inventory--search
comparison. The unresolved question is how visual input affects the coverage
and stability of measurements drawn from a restricted response set, and what
those properties allow an auditor to infer from the resulting scores.

This matters because a conditional preference between selected answers can
remain numerically plausible even when those answers capture little of the
output distribution. A rating can also have high response coverage while
changing under an equivalent scale definition. Coverage and robustness
therefore address different requirements for interpreting the measurement;
neither can be assumed from the choice to read probabilities alone.

The present study addresses this gap through two linked evaluations.
The inventory study addresses RQ1 by examining gender-related responses
while separating general agreement, item discrimination and response coverage.
The search study addresses RQ2 by examining rating robustness and comparing
model judgements with the historical human reference. Together, these evaluations assess the reliability
of the measurements before considering whether inventory responses correspond
with downstream judgements. The gap is defined relative to the reviewed
approaches, without claiming that every related method has been excluded.

% \begin{plan}
% One page. Summarise as \cref{tab:gap}, then one paragraph: prior work measures
% either generated text or embeddings; probability-level work is text-only;
% human-anchored work is either text-only or human-only; none reports how much
% of the distribution the measurement captured. \textbf{Verify every cell
% against the paper while writing}; cells marked ``?'' are not yet checked.
% \textbf{Done (Nov 2026):} every cell checked against the paper's abstract and
% reported method; no ``?'' remains.
% \end{plan}
